% !TEX root = ../Thesis.tex
\chapter{Abstract}

% Revise the final paragraph once the SnapQuery comparison and Phase 2 are done.

Clinical knowledge graphs store medical facts as entities and the relations
between them: a diagnosis, the drug given for it, the laboratory value that
prompted a change of dose. Every fact can be checked on its own, which is what
makes such a graph appealing as the backend for a question answering system. The
difficulty is size. A hospital graph serialises to far more text than a language
model can read, so before any question can be answered, something has to decide
which fragment of the graph the model will be shown.

This thesis studies that decision on the Swiss Personalized Health Network (SPHN)
graph of the Swiss Transplant Cohort Study at University Hospital Basel. It
compares the Prize-Collecting Steiner Tree (PCST) formulation introduced by
G-Retriever with SnapQuery, the assistant that clinicians and data managers at
the STCS Data Center already rely on. Both are scored on answer-node recall at a
given subgraph size: how much of a correct answer the retrieved fragment
contains, relative to how large that fragment is. Two ablations, similarity-only
top-$k$ selection and breadth-first expansion from seed nodes, separate what the
connectivity constraint contributes from what the similarity scores contribute.

Measurement showed the graph to differ from the benchmarks G-Retriever was
validated on in ways that change how the method behaves. It holds 21{,}537{,}305
nodes for 1{,}197 patients. A single reified laboratory construct accounts for
95.3\% of them. A few provenance nodes carry a large share of all edges, one of
them fifteen per cent by itself. Text is scarce and heavily shared, so that one
patient's 15{,}810 prepared nodes rest on only 614 distinct descriptions.

Three consequences follow from G-Retriever's prize and cost definitions, and all
three were then confirmed by experiment. The edge-prize mechanism dilutes to
nothing when relation types repeat at high multiplicity, and under the published
parameters the method returned the entire graph it was given. Node selection is
settled by tie-breaking rather than by similarity, so a score can identify which
kind of entity a question concerns but not which instance of it. Connectivity
costs almost nothing to satisfy once high-degree provenance nodes are present,
and at matched subgraph size every strategy tested recalled the same amount,
whether it enforced connectivity or ignored it entirely.

Two further results came out of the same measurements. Resolving terminology
codes to canonical descriptions, instead of embedding whatever label the graph
happens to carry, lifts diagnoses that no wording could retrieve from zero recall
to 0.80. Every retrieval strategy gains equally, which places the step in graph
preparation rather than in retrieval, and the gain reverses completely when the
catalogue is written in a different language from the question. A
candidate-generation stage brings the method within reach at patient scale
without costing recall, provided terminology resolution has already been applied.

The comparison against SnapQuery has not yet been run. All retrieval experiments
reported here use a deterministic lexical encoder on a synthetic graph built to
reproduce the measured structure of one real patient, so they describe how the
method behaves on this class of graph rather than how well it retrieves clinical
answers.
